% Figure: the Coriolis deflection seen on a rotating turntable. A puck
% slid radially outward along a straight line in the inertial frame
% appears, to an observer co-rotating with the turntable, to curve to
% one side. The two paths are drawn: the straight inertial path and the
% curved path painted on the turntable.
\begin{figure}[htbp]
\centering
\begin{tikzpicture}[lbl/.style={font=\small}]
% turntable disc
\draw[enggray, very thick] (0,0) circle (3.2);
\fill[enggray!12] (0,0) circle (3.2);
\fill[engblack] (0,0) circle (2pt);
\node[lbl, below left] at (0,0) {centre};
% rotation arrow (counter-clockwise)
\draw[-{Stealth[length=2.5mm]}, engblue, thick] (135:3.5) arc (135:95:3.5);
\node[engblue, lbl] at (118:3.85) {$\Omega$};
% inertial straight path: launched at center toward angle 0 (to the right)
\draw[engred, very thick, -{Stealth[length=3mm]}] (0,0) -- (3.0,0);
\node[engred, lbl, above] at (2.2,0.05) {inertial path (straight)};
% curved path as seen on the rotating disc: deflects to the right of
% travel for counter-clockwise rotation. Plot a curve bending downward.
\draw[englime!60!black, very thick, -{Stealth[length=3mm]}]
  (0,0) .. controls (1.4,-0.35) and (2.4,-1.1) .. (2.7,-1.55);
\node[englime!50!black, lbl, below] at (2.0,-1.4) {path on the disc (curved)};
% small deflection bracket
\draw[enggray, dashed] (3.0,0) -- (2.7,-1.55);
\node[enggray, lbl, right] at (3.0,-0.8) {deflection};
\end{tikzpicture}
\caption[Coriolis deflection on a rotating turntable]{A puck launched
from the centre of a counter-clockwise turntable travels in a straight
line in the inertial frame (the red arrow), because once it leaves the
centre no horizontal force acts on it. An observer painting the puck's
track on the rotating disc records a curved path (the green arrow) that
bends to the right of the direction of travel. The curvature is the
Coriolis effect: in the rotating frame the puck obeys
$m\vec{a}'=-2m\vec{\Omega}\times\vec{v}'$, a deflecting pseudo-force
perpendicular to the puck's velocity. The same deflection sets the
rotation sense of large-scale winds and biases any projectile that
flies far enough for the planet to turn beneath it.}
\label{fig:vol03ch04:coriolis}
\end{figure}
